\documentclass[exercice,a4]{thedocument}%
\usepackage[french]{babel}%
\usepackage{textcomp}%
\usepackage{pstricks}%
\usepackage{scratch3}%
\usepackage{tabularx}%
\usepackage{enumitem}%
\startdatakeys
    \source{21GENMATNC1}
    \BO{2026}
    \cycle{4}
    \competencesDuSocle{CA,RA}
    \objectifsApprentissage{C4B3ca,C4D2ba}
    \motsCles{Thalès,droites parallèles,trigonométrie,angle,Scratch,programme}
\enddatakeys
\begin{document}%
\def\degres{\textsuperscript{\textdegree}}%

\medskip

On dispose des informations suivantes sur le phare Amédée, une balise et une bouée :

\begin{itemize}
\item la hauteur du phare est de $56$~m ;
\item la balise est située à $260$~m du phare;
\item la balise et la bouée sont distantes de $151$~m ;
\item la bouée O, le sommet B de la balise et le sommet P du phare sont considérés comme trois points alignés.
\end{itemize}

Schéma de la situation :

\begin{center}
\psset{unit=1cm,arrowsize=2pt 3}
\begin{pspicture}(10.5,3.5)
\psline(0,0.5)(10,0.5)(0,3.1)%HOP
\psline[linestyle=dashed](1.3,0.5)(1.3,2.7)
\psline[linestyle=dashed](7.1,0.5)(7.1,1.25)
\uput[u](1.3,2.7){P} \uput[d](1.3,0.5){H} \uput[u](7.1,1.25){B} \uput[d](7.1,0.5){A}
\uput[d](10,0.5){O} \uput[l](1.3,1.6){56 m} \uput[d](4.2,0.5){260 m} \uput[d](8.55,0.5){151 m}
\rput{90}(1.6,1.6){Phare} \rput{90}(7.4,0.85){Balise}\rput{90}(10,1){Bouée}
\psline[linewidth=0.6pt]{<->}(1.3,0)(7.1,0)
\psline[linewidth=0.6pt]{<->}(10,0)(7.1,0)
\end{pspicture}
\end{center}

\textbf{Les droites (PH) et (BA) sont parallèles.}

\medskip

\begin{enumerate}
\item Quelle est la distance OH en m ?
\item Déterminer la hauteur AB de la balise. Arrondir au dixième de m près.

\textbf{Rédiger la réponse en faisant apparaître les différentes étapes.}
\end{enumerate}

Le haut du phare est protégé par une barrière composée de sculptures.

\begin{center}
Contour de la sculpture
\medskip

\psset{unit=1cm,arrowsize=2pt 3}
\begin{pspicture}(-1.5,-1.5)(1.5,1.5)
\pspolygon(0.9;0)(1.4;15)(0.9;30)(1.1;45)(0.9;60)(1.4;75)(0.9;90)(1.1;105)(0.9;120)(1.4;135)(0.9;150)(1.1;165)(0.9;180)(1.4;195)(0.9;210)(1.1;225)(0.9;240)(1.4;255)(0.9;270)(1.1;285)(0.9;300)(1.4;315)(0.9;330)(1.1;345)
\end{pspicture}
\end{center}

On souhaite réaliser un programme Scratch pour reproduire le contour de cette sculpture.
\begin{enumerate}[resume]
\item Calculer la mesure de l'angle $\widehat{a}$ en degré dans la figure ci-dessous:

\begin{center}
\psset{unit=1cm,arrowsize=2pt 3}
\begin{pspicture}(-1.5,-1.5)(1.5,1.5)
\psline[linestyle=dotted,linewidth=1.5pt](-1.5,0)(1.5,0)
\psline[linestyle=dotted,linewidth=1.5pt](0,0)(1.5;-72)
\psarc(0,0){0.5}{298}{360}\psarc(0,0){0.6}{0}{180}\psarc(0,0){0.5}{180}{288}\psarc(0,0){0.45}{180}{288}
\rput(0.9;90){$180\degres$}\rput(0.8;334){$72\degres$}\rput(0.75;234){$\widehat{a}$}
\psline(-0.6,0)(0,0)(0.5;288)
\end{pspicture}
\end{center}

Le script 1 permet de tracer le motif en pointillé ci-dessous (on part du point A et on s'arrête au point B).

\begin{center}
\psset{unit=1cm,arrowsize=2pt 3}
\begin{pspicture}(4.6,6)
%\psgrid
\psline(2.5,0)(0.2,0.3)(2.2,1.9)(1.1,3.2)(2.7,3.4)(2.4,5.5)(4.4,4.2)
\uput[l](0.2,0.3){B}\uput[r](0.4,0.45){$36\degres$}
\uput[r](1.2,3.1){$72\degres$}\uput[u](1.8,3.3){28}\uput[dl](1.6,2.6){28}
\uput[l](2.5,4.5){38}\uput[ul](2.2,5.5){A}\uput[ur](3.4,4.8){38}\uput[ul](1.2,1.18){38}
\uput[dl](1.4,0.1){38}
\rput{126}(2.2,1.9){\psframe(0.2,0.2)}
\rput{98}(2.7,3.4){\psframe(0.2,0.2)}
\end{pspicture}
\end{center}

\item Compléter le script 1 de l'annexe.
\end{enumerate}

Le script final permet de réaliser le contour de la sculpture.

\begin{enumerate}[resume]
\item Compléter le script final de l'annexe.
\end{enumerate}

\begin{center}
\begin{tabularx}{\linewidth}{XX}
\textbf{Script 1} &\textbf{Script final}\\
\begin{scratch}
\initmoreblocks{définir \namemoreblocks{motif}}%
\blockmove{avancer de \ovalnum{38} pas}
\blockmove{tourner \turnright{} de \ovalnum{} degrés}
\blockmove{avancer de \ovalnum{28} pas}
\blockmove{tourner \turnleft{} de \ovalnum{108} degrés}
\blockmove{avancer de \ovalnum{} pas}
\blockmove{tourner \turnright{} de \ovalnum{90} degrés}
\blockmove{avancer de \ovalnum{} pas}
\end{scratch}&
\begin{scratch}
\blockinit{Quand \greenflag est cliqué}
\blockpen{effacer tout}
\blockpen{stylo en position d'écriture}
\blockrepeat{répéter \ovalnum{} fois}
{
\blocklook{motif}
\blockmove{tourner \turnleft{} de \ovalnum{90} degrés}
}
\end{scratch}
\end{tabularx}
\end{center}

\showthesource
\end{document}
